\section{Generated submodules}
\label{sec:generated}
Throughout this section, $K$ is a field and $X=X(\vec s;K)$.
We describe the submodule $\langle T\rangle$ from Section~2
by its normalized kernel. This gives finite-dimensional data for
computing its homology.

Let $P=K[\Delta^n]$ be the free simplicial $K$-module on the standard
$n$-simplex, so $P_p$ has basis the nondecreasing maps $[p]\to[n]$.
For $T\in X_n$, its representing property gives
\[
 f_T:P_\bullet\longrightarrow X_\bullet,
 \qquad f_T(\alpha)=X(\alpha)(T)
 \quad\bigl(\alpha:[p]\to[n]\bigr).
\]
Its image is the simplicial submodule $\langle T\rangle$ generated by $T$,
and $P/\ker f_T\cong\langle T\rangle$. Put
\[
 C=N(P),\qquad g_T=N(f_T),\qquad L(T)=\ker g_T.
\]
Normalization is exact, so
\begin{equation}\label{eq:generated-normalization}
 N(\langle T\rangle)\cong C/L(T).
\end{equation}
Here $C_p$ has dimension $\binom{n+1}{p+1}$ for $0\leq p\leq n$
and is zero otherwise. Hence the finitely many spaces $L_p(T)$ determine
the entire simplicial kernel by Dold--Kan. They classify the quotient
with its distinguished generating $n$-simplex; they need not determine
its inclusion in $X$, and different kernels can yield isomorphic
simplicial modules after the generator is forgotten.

Homology is computed from normalized complexes. Since $H_0(C)=K$ and
$H_j(C)=0$ for $j>0$, the short exact sequence
$0\to L(T)\to C\to C/L(T)\to0$ gives
\begin{equation}\label{eq:generated-homology}
 \begin{gathered}
 H_j(\langle T\rangle)\cong H_{j-1}(L(T))\qquad(j\geq2),\\
 0\longrightarrow H_1(\langle T\rangle)\longrightarrow H_0(L(T))
 \xrightarrow{\epsilon}K\longrightarrow H_0(\langle T\rangle)
 \longrightarrow0,
 \end{gathered}
\end{equation}
where $\epsilon$ is induced by the augmentation sending every vertex of
$\Delta^n$ to $1$.

For the matrix in Example~\ref{ex:matrix-cycle},
$N(\langle E_{01}\rangle)$ is $K$ in degree $1$ and zero elsewhere:
all its higher simplices are degenerate.

\appendix
\section{Kernel families over an infinite field}
\label{app:kernel-families}
Assume now that $K$ is infinite. Retain the notation $P$, $C$, $g_T$,
and $L(T)$ from Section~\ref{sec:generated}. We parameterize the kernels
on the locus where their dimensions are minimal.

For a finite-dimensional $K$-vector space $V$ and $0\leq r\leq\dim V$,
write $\operatorname{Gr}(r,V)$ for the Grassmannian of $r$-dimensional
linear subspaces of $V$. This is a projective $K$-variety whose points
over a field extension $F/K$ are the $r$-dimensional $F$-linear subspaces
of $V\otimes_K F$. Its tautological subbundle of the trivial bundle
with fiber $V$ has fiber $L$ at the point representing $L\subseteq V$.

View $W=X_n$ as an affine space. For $0\leq p\leq n$, let $r_p$ be
the maximal rank of $g_{T,p}$ and put $\ell_p=\dim C_p-r_p$.
The locus $U\subset W$ on which all these ranks are attained is a
nonempty open subset. Constant-rank kernels define a morphism
\begin{equation}\label{eq:kernel-map}
 \Psi:U\longrightarrow G=\prod_{p=0}^n\operatorname{Gr}(\ell_p,C_p),
 \qquad T\longmapsto(L_p(T))_p.
\end{equation}
Its values satisfy $d_0L_p\subseteq L_{p-1}$, so they are chain
subcomplexes of $C$. By Chevalley's theorem \cite[Tag 054K]{Stacks},
the image $\mathcal M=\Psi(U)$ is constructible. Give its Zariski
closure $Y=\overline{\mathcal M}$ the reduced
structure. When $\mathcal M$ is locally closed, it too is given the
reduced induced structure. Thus $\mathcal M$
records marked quotients of $P$, rather than isomorphism classes of
unmarked simplicial modules. A $K$-variety is \emph{rational over $K$}
if it is birational over $K$ to an affine space $\mathbb A_K^d$.

\begin{proposition}\label{prop:kernel-rationality}
The variety $\overline{\mathcal M}$ is rational over $K$.
\end{proposition}
\begin{proof}
The closure is integral because $U$ is irreducible. Each $g_{T,p}$
depends linearly on $T$. For a tuple $L=(L_p)$, the conditions
\[
 g_{T,p}(L_p)=0\quad(0\leq p\leq n)
\]
therefore define a linear subspace $V_L\subset W$.
Let $\mathcal E\subset W\times Y$ be the incidence scheme cut out by
these conditions on the tautological subbundles $L_p$ pulled back from
the factors of $G$. Over $U\times Y$,
$L_p$ and $\ker g_{T,p}$ are subbundles of equal rank, so inclusion
forces equality. Hence $\mathcal E\cap(U\times Y)$ is the graph of
$\Psi$ as a scheme, and its fiber over $L$ is $U\cap V_L$.

Let $F=K(Y)$ be the function field of $Y$, let $\eta$ be its generic
point, and write $W_F,U_F$ for base change from $K$ to $F$.
Let $V_\eta\subset W_F$ be the linear subspace defined by the incidence
conditions at $\eta$. Its intersection with $U_F$
is nonempty, since $\Psi$ is dominant onto its image closure. Put
$d=\dim_F V_\eta$. Choose $d$ coordinate functions on $W$ whose
restriction gives an isomorphism $V_\eta\to F^d$, and denote their
joint map by $\pi:W\to K^d$. Solving the defining linear equations
expresses the unique element of $V_\eta$ with $\pi(T)=c$ rationally
in the kernel coordinates and linearly in $c$. The condition that this
element belongs to $U_F$ defines a nonempty open subset of $F^d$.
Since $K$ is infinite, $K^d$ is Zariski dense in $F^d$; choose
$c\in K^d$ in that open subset.

After restricting to a nonempty open subset of
$\overline{\mathcal M}$, the subspaces $V_L$ have constant dimension,
$\pi|_{V_L}$ is invertible, and the solutions with $\pi(T)=c$
belong to $U$. Every fiber there consequently meets the affine slice
$S=\pi^{-1}(c)$ in exactly one point, varying rationally with $L$.
These solutions and $\Psi|_S$ are inverse on nonempty open subsets.
As $S\cong\mathbb A_K^{\dim W-d}$, the closure is rational.
\end{proof}

For $V=K^3$, a plane is the kernel of a nonzero linear form $(a,b,c)$,
unique up to a nonzero scalar. Thus
\[
 \operatorname{Gr}(2,K^3)\xrightarrow{\ \sim\ }\mathbb P_K^2,
 \qquad \ker(a,b,c)\longmapsto[a:b:c],
\]
where $[a:b:c]$ denotes the corresponding homogeneous coordinates.

\begin{proposition}\label{prop:three-by-three}
For shape $(3,3)$, the generic normalized kernel dimensions are
$(\ell_0,\ell_1,\ell_2)=(2,0,0)$, and the kernel image is
\[
 \mathcal M\cong\mathbb P_K^2\setminus\{[1:1:1]\}.
\]
For every $T\in U$,
\[
 H_1(\langle T\rangle)\cong K,
 \qquad H_j(\langle T\rangle)=0\quad(j\ne1).
\]
\end{proposition}
\begin{proof}
Write $T=(t_{ij})$ and $(a,b,c)=(t_{00},t_{11},t_{22})$.
In degree zero, $f_T$ is the row $(a,b,c):K^3\to K$.
In degree one, the columns corresponding to the three constant maps
$[1]\to[2]$ are $aJ,bJ,cJ$, where $J$ is the all-ones $2\times2$
matrix. When $(a,b,c)\ne0$, they span $KJ$. Modulo this line, the
other three columns are represented by
\[
 B_T=\begin{pmatrix}
 t_{01}-a&t_{02}-a&t_{12}-b\\
 t_{10}-a&t_{20}-a&t_{21}-b\\
 b-a&c-a&c-b
 \end{pmatrix}.
\]
Its first two rows are arbitrary, while its last row is nonzero
exactly when $a,b,c$ are not all equal. Consequently $f_{T,1}$ can
have rank $4$ for a specified diagonal precisely when that diagonal
is nonconstant. Such a choice gives $\rank g_{T,0}=1$ and
$\rank g_{T,1}=3$, using the natural decomposition
$P_1=C_1\oplus s_0C_0$ and its counterpart for $X$.

The normalized simplex complex $C=N(K[\Delta^2])$ has dimensions
$(3,3,1)$ and injective differential $C_2\to C_1$. Since $g_{T,1}$
is injective, the chain-map identity makes $g_{T,2}$ injective as
well. These are maximal ranks in all degrees, and
$U=\{T:\det B_T\ne0\}$. Thus on $U$ the kernel
complex is $\ker(a,b,c)$ in degree zero and zero elsewhere.
Conversely, every nonconstant projective diagonal is realized by
choosing the first two rows of the displayed matrix to complete its
last row to a basis. This construction works over every residue field.
Under the displayed identification of $\operatorname{Gr}(2,K^3)$ with
$\mathbb P_K^2$, the kernel map therefore has exactly the stated image.

The restriction of the augmentation $(1,1,1)$ to $\ker(a,b,c)$ is
nonzero, because $(a,b,c)$ is not proportional to $(1,1,1)$. It is
therefore surjective with one-dimensional kernel. Formula
\eqref{eq:generated-homology} proves the homology assertions.
\end{proof}

For comparison, the all-ones $3\times3$ tensor generates the constant
simplicial module $K$: every simplicial operator sends it to an
all-ones tensor. Its homology is $K$ in degree zero and zero in
positive degrees. Thus generated submodules distinguish tensors by
homology even when their common ambient module is contractible.
